\chapter{Gateway (Core Engine)}

\section{Overview}

The Gateway is the core process of OpenClaw, responsible for:

\begin{itemize}
  \item \textbf{Message Routing}: Receiving inbound messages from channels and routing them to the correct Agent
  \item \textbf{Session Management}: Maintaining all session state (single source of truth)
  \item \textbf{Agent Runtime}: Embedding pi-agent-core to execute the agent loop
  \item \textbf{RPC Service}: Exposing WebSocket/HTTP APIs
  \item \textbf{Plugin Host}: Loading and running all registered plugins
\end{itemize}

Config path: \texttt{gateway.*}

\section{Startup Methods}

\subsection{Foreground Mode}

\begin{lstlisting}[language=bash]
openclaw gateway --port 18789
\end{lstlisting}

\subsection{Daemon Mode}

\begin{lstlisting}[language=bash]
openclaw daemon install   # Install systemd/launchd service
openclaw daemon start     # Start the daemon
openclaw daemon status    # Check status
\end{lstlisting}

\section{Agent Loop}

Each inbound message triggers a complete agent loop:

\begin{center}
\begin{tikzpicture}[
    node distance=0.6cm,
    agentstep/.style={draw=teal!60, fill=teal!8, rounded corners=4pt, minimum width=12cm, minimum height=0.7cm, font=\small, align=left},
    arr/.style={-{Stealth[length=2.5mm]}, thick, draw=teal!60}
]
    \node[agentstep] (s1) {1. Receive inbound message, validate params, resolve session};
    \node[agentstep, below=of s1] (s2) {2. Route to target Agent, load Skills snapshot};
    \node[agentstep, below=of s2] (s3) {3. Assemble system prompt (base + Skills + bootstrap context)};
    \node[agentstep, below=of s3] (s4) {4. Call model for inference, stream assistant response};
    \node[agentstep, below=of s4] (s5) {5. Execute tool calls (possibly in sandbox)};
    \node[agentstep, below=of s5] (s6) {6. Persist session, send reply to channel};

    \draw[arr] (s1) -- (s2);
    \draw[arr] (s2) -- (s3);
    \draw[arr] (s3) -- (s4);
    \draw[arr] (s4) -- (s5);
    \draw[arr] (s5) -- (s6);
\end{tikzpicture}
\end{center}

\section{RPC Interface}

The Gateway exposes RPC interfaces via WebSocket and HTTP:

\begin{itemize}
  \item \texttt{agent} / \texttt{agent.wait}: Trigger agent run
  \item \texttt{sessions.list}: List sessions
  \item \texttt{voicecall.*}: Voice call (plugin)
  \item \texttt{message.send}: Send message to channel
\end{itemize}

Example (calling RPC via CLI):

\begin{lstlisting}[language=bash]
openclaw gateway call sessions.list --params '{}'
\end{lstlisting}

\section{Authentication}

\textbf{Anthropic API key} is recommended:

\begin{lstlisting}[language=bash]
export ANTHROPIC_API_KEY="..."
# Or store in daemon env file:
echo 'ANTHROPIC_API_KEY=...' >> ~/.openclaw/.env
\end{lstlisting}

OAuth and setup-token flows are also supported. Use \texttt{openclaw models status} to verify authentication status.

\section{Common Operations}

\begin{lstlisting}[language=bash]
openclaw doctor            # Health check
openclaw status            # Show storage paths & recent sessions
openclaw models status     # Check model auth
openclaw gateway --port N  # Start with custom port
\end{lstlisting}
